\section*{第 \ref{ch:b1:weathering-clays} 章——风化与黏土}

\begin{solution}{exo:b1:weathering-clays:1}
\ce{K+}：$1/1.38 = 0.72$；\ce{Fe^2+}：$2/0.78 = 2.56$；均小于 3，为可溶
阳离子。\ce{Fe^3+}：$3/0.645 = 4.65$；\ce{Ti^4+}：$4/0.605 = 6.61$；均
介于 3 与 10 之间，为\omterm{def:b1:weathering-clays:hydrolysis}{水解}阳离子。在缺氧的排水不良土壤中，铁以
\ce{Fe^2+} 形式释放，留在溶液中，被水带走；在排水良好的土壤中，
氧气将其氧化为 \ce{Fe^3+}，后者立即以氧化物和氢氧化物的形式沉淀，
留存于原地。
\end{solution}

\begin{solution}{exo:b1:weathering-clays:2}
A 在乙二醇中膨胀，加热后塌缩，为\omterm{def:b1:weathering-clays:clay-mineral}{蒙皂石}。B 位于
\SI{7}{\angstrom} 处，在 \SI{550}{\celsius} 时被破坏，为\omterm{def:b1:weathering-clays:clay-mineral}{高岭石}。C 位于
\SI{10}{\angstrom} 处，无论经何种处理，为\omterm{def:b1:weathering-clays:clay-mineral}{伊利石}。D 位于 \SI{14}{\angstrom}
处，不膨胀也不塌缩，为\omterm{def:b1:weathering-clays:clay-mineral}{绿泥石}。
\end{solution}

\begin{solution}{exo:b1:weathering-clays:3}
两个正\omterm{def:b1:rock-forming-minerals:groups}{长石}提供一个\omterm{def:b1:weathering-clays:clay-mineral}{高岭石}所需的 2 个铝；$6 - 2 = 4$ 个硅
以 \ce{H4SiO4} 形式离开；2 个 \ce{K+} 离开，由 2 个 \ce{H+} 平衡；
氧原子平衡 $16 + w = 9 + 16$，得 $w = 9$：
\[
  \ce{2 KAlSi3O8 + 2 H+ + 9 H2O -> Al2Si2O5(OH)4 + 2 K+ + 4 H4SiO4}.
\]
氢原子核对：$2 + 18 = 20 = 4 + 16$。钾全部离开，硅离开的比例为 $4/6 =
2/3$。
\end{solution}

\begin{solution}{exo:b1:weathering-clays:4}
(a) 质量不变：$10 \times 999.87/916.72 = \qty{10.91}{cm^3}$。
(b) 约在 \qty{-4}{\celsius} 至 \qty{-15}{\celsius} 之间。刚低于零度时，
冰无法积累足够的压强撑开裂缝；在更严寒的条件下，
水不再通过孔隙向冰迁移，冰
停止增长。(c) 由
\cref{prop:b1:weathering-clays:surface}可知，取 $n = \qty{1}{m}/\qty{1}{cm} = 100$，
面积乘以 100，从 \qty{6}{m^2} 增至 \qty{600}{m^2}。
\end{solution}

\begin{solution}{exo:b1:weathering-clays:5}
$S/m = 6/(\rho d)$：当 \qty{1}{cm} 时，$6/(2600\times0.01) =
\qty{0.231}{m^2.kg^{-1}}$；当 \qty{1}{mm} 时，\qty{2.31}{m^2.kg^{-1}}；当
\qty{1}{\micro m} 时，\qty{2308}{m^2.kg^{-1}}，即 \qty{2.3}{m^2.g^{-1}}。
化学反应发生在颗粒表面，因此其每千克岩石的反应速率
随比表面积增大，从厘米到微米增大了一万倍。
\omterm{def:b1:weathering-clays:weathering}{物理风化}将岩石切割成越来越小的碎块，
并张开裂隙使水分渗入：它成倍增加了
\omterm{def:b1:weathering-clays:weathering}{化学风化}得以作用的表面。
\end{solution}

\begin{solution}{exo:b1:weathering-clays:6}
一吨钠\omterm{def:b1:rock-forming-minerals:groups}{长石}为 $\num{e6}/262.22 = \qty{3813.6}{mol}$。由
\cref{prop:b1:weathering-clays:feldspars}可知，两摩尔产生一摩尔\omterm{def:b1:weathering-clays:clay-mineral}{高岭石}：
$\qty{1906.8}{mol} \times \qty{258.16}{g.mol^{-1}} = \qty{492.3}{kg}$。三个硅原子中的两
个离开，形成 $2\times3813.6 = \qty{7627.2}{mol}$ 的
\ce{SiO2}，\qty{458.2}{kg}。钠全部离开：$3813.6\times22.99 =
\qty{87.7}{kg}$。留存为黏土的质量不足\omterm{def:b1:rock-forming-minerals:groups}{长石}的一半。
\end{solution}

\begin{solution}{exo:b1:weathering-clays:7}
由\cref{prop:b1:weathering-clays:immobile}可知，每千克花岗岩留下
$0.30/0.60 = \qty{0.5}{kg}$ 的腐岩。损失分数：钠为 $1 -
(0.35/3.5)(0.5) = 0.95$；硅为 $1 - (60/72)(0.5) = 0.583$；铝为 $1 -
(27/14)(0.5) = 0.036$。铝基本保留，钠几乎
全部流失，硅也流失了近 \qty{60}{\percent}，尽管残留物的
大部分保存于石英中：\omterm{def:b1:rock-forming-minerals:groups}{长石}正在转化为\omterm{def:b1:weathering-clays:clay-mineral}{高岭石}，
即\omterm{def:b1:weathering-clays:hydrolysis-degree}{单硅铝化}。
\end{solution}

\begin{solution}{exo:b1:weathering-clays:8}
叶蜡石：$x = y = 0$，电荷为 0。白\omterm{def:b1:rock-forming-minerals:groups}{云母}层：$x = 1$，电荷为 $-1$。
贝得石层：$x = 0.5$，电荷为 $-0.5$。蒙脱石层：$y = 0.33$，
电荷为 $-0.33$。两个\omterm{def:b1:weathering-clays:clay-mineral}{蒙皂石}层，即贝得石层和蒙脱石层，能
膨胀：它们电荷较小，由数量较少的层间阳离子平衡，这些阳离子
松散地连接各层，并在水化时吸入水分。叶蜡石没有
可水化的层间阳离子，白\omterm{def:b1:rock-forming-minerals:groups}{云母}的电荷 $-1$ 由能锁住各层的
钾离子平衡。
\end{solution}

\begin{solution}{exo:b1:weathering-clays:9}
(a) Mg：2 与 2；Si：1 与 1；C：4 与 4；O：$4 + 8 + 4 = 16$ 与 $12 + 4 =
16$；H：8 与 $4 + 4 = 8$；电荷：0 与 $4 - 4 = 0$。(b) 镁\omterm{def:b1:rock-forming-minerals:groups}{橄榄石}的
摩尔质量为 $2\times24.305 + 28.085 + 4\times15.999 =
\qty{140.69}{g.mol^{-1}}$，因此一吨即 \qty{7108}{mol}。其消耗
$4\times7108 = \qty{28431}{mol}$ 的 \ce{CO2}，\qty{1251}{kg}；每个化学式单位含两个
镁，由\cref{prop:b1:weathering-clays:co2}可知，储存
$2\times7108$ 摩尔，即 \qty{626}{kg}，为消耗量的一半。(c) 一吨
方解石，即 \qty{9991}{mol}，消耗 \qty{9991}{mol} 的 \ce{CO2}（\qty{440}{kg}）
溶解时吸收，并在碳酸盐重新沉淀时全部释放：
因此没有任何储存。
\end{solution}

\begin{solution}{exo:b1:weathering-clays:10}
(a) 由 $3 < 3/r < 10$ 得 $\qty{0.3}{\angstrom} < r < \qty{1.0}{\angstrom}$。
(b) $3/0.27 = 11.1 > 10$：硼以含氧阴离子形式迁移，即硼酸根。(c) $2/0.83 =
2.41$，为可溶阳离子；$4/0.53 = 7.55$，为\omterm{def:b1:weathering-clays:hydrolysis}{水解}阳离子。在贫氧的水中，
锰以 \ce{Mn^2+} 形式溶解；泉水与空气接触处，
它被氧化为 \ce{Mn^4+}，并立即以黑色氧化物的形式沉淀在
岩石表面。(d) 其\omterm{def:b1:weathering-clays:ionic-potential}{离子势} $4/0.40 = 10.0$，位于
\omterm{def:b1:weathering-clays:hydrolysis}{水解}阳离子与含氧阴离子的分界线上：硅既不像铝那样被保留，
也不像硫那样被带走，浸出比例取决于
流经的水量，这正是三级\omterm{def:b1:weathering-clays:hydrolysis}{水解}程度的依据。
\end{solution}

\begin{solution}{exo:b1:weathering-clays:11}
(a) 以碳的摩尔质量 \qty{12.011}{g.mol^{-1}} 计，\num{11.7e12} 至
\qty{13.6e12}{mol} 即每年 0.141 至 0.163 十亿吨碳；以
\qty{44.009}{g.mol^{-1}} 计，为 0.515 至 0.599 十亿吨二氧化碳。
(b) $4.08/13.6 = \qty{30}{\percent}$ 至 $4.08/11.7 = \qty{35}{\percent}$。
(c) 消耗的二氧化碳中有一半，在钙和镁沉淀为碳酸盐时被储存，
即每年 0.070 至 0.082 十亿吨碳。实际数值
更小，因为部分碳酸氢根随钠和
钾一起迁移，而钠、钾在海底不形成碳酸盐。
\end{solution}

\begin{solution}{exo:b1:weathering-clays:12}
Si/Al $= (w_{\ce{SiO2}}/60.08)/(2w_{\ce{Al2O3}}/101.96)$. A: $0.7748/0.7748 =
1.00$，\omterm{def:b1:weathering-clays:clay-mineral}{高岭石}的比值，\omterm{def:b1:weathering-clays:hydrolysis-degree}{单硅铝化}。B: $0.9987/0.4904 = 2.04$，
2:1 型\omterm{def:b1:weathering-clays:clay-mineral}{黏土矿物}，\omterm{def:b1:weathering-clays:hydrolysis-degree}{双硅铝化}。C: $0.0333/1.2358 = 0.027$，几乎不
剩硅，三水铝石，\omterm{def:b1:weathering-clays:hydrolysis-degree}{铝铁化}。从淋滤最弱到最强，
即朝向最热、最湿的气候：B、A、C。
\end{solution}

\begin{solution}{pb:b1:weathering-clays:1}
\textbf{1.} \ce{2 NaAlSi3O8 + 2 H+ + 9 H2O -> Al2Si2O5(OH)4 + 2 Na+ + 4 H4SiO4}.

\textbf{2.} \ce{CaAl2Si2O8 + 2 H+ + H2O -> Al2Si2O5(OH)4 + Ca^2+}。没有
硅析出：钙\omterm{def:b1:rock-forming-minerals:groups}{长石}本身已具有\omterm{def:b1:weathering-clays:clay-mineral}{高岭石}中每个铝对应的一个硅。

\textbf{3.} 一个三水铝石吸收一个铝，相当于半个钙\omterm{def:b1:rock-forming-minerals:groups}{长石}，加倍后得：
\ce{CaAl2Si2O8 + 2 H+ + 6 H2O -> 2 Al(OH)3 + Ca^2+ + 2 H4SiO4}。氧：
$8 + 6 = 6 + 8$；氢：$2 + 12 = 6 + 8$。

\textbf{4.} \ce{5 NaAlSi3O8 + 4 H+ + 16 H2O -> 2 Na_{0.5}Al2(Si_{3.5}Al_{0.5})O10(OH)2 + 4 Na+ + 8 H4SiO4}; % ce-unbalanced-ok: decimal subscripts the checker cannot read; balance tested in tests/bachelor-1/test_weathering_clays.py
氢：左侧 $4 + 32 = 36$，右侧 $2\times2 + 8\times4 = 36$，
两侧相等，配平。

\textbf{5.} 每个 \ce{H+} 伴随一个 \ce{HCO3-}，它来自一个 \ce{CO2} 和一个
\ce{H2O}:
\ce{2 NaAlSi3O8 + 2 CO2 + 11 H2O -> Al2Si2O5(OH)4 + 2 Na+ + 2 HCO3- + 4 H4SiO4}
和 \ce{CaAl2Si2O8 + 2 CO2 + 3 H2O -> Al2Si2O5(OH)4 + Ca^2+ + 2 HCO3-}。

\textbf{6.} 钠\omterm{def:b1:rock-forming-minerals:groups}{长石}生成贝得石为\omterm{def:b1:weathering-clays:hydrolysis-degree}{双硅铝化}；钠\omterm{def:b1:rock-forming-minerals:groups}{长石}与钙\omterm{def:b1:rock-forming-minerals:groups}{长石}生成
\omterm{def:b1:weathering-clays:clay-mineral}{高岭石}为\omterm{def:b1:weathering-clays:hydrolysis-degree}{单硅铝化}；钙\omterm{def:b1:rock-forming-minerals:groups}{长石}生成三水铝石为\omterm{def:b1:weathering-clays:hydrolysis-degree}{铝铁化}。

\textbf{7.} 钠\omterm{def:b1:rock-forming-minerals:groups}{长石}：$\num{e6}/262.22 = \qty{3813.6}{mol}$；
钙\omterm{def:b1:rock-forming-minerals:groups}{长石}：$\num{e6}/278.20 = \qty{3594.5}{mol}$。

\textbf{8.} 钠\omterm{def:b1:rock-forming-minerals:groups}{长石}：$3813.6/2 = \qty{1906.8}{mol}$ 的\omterm{def:b1:weathering-clays:clay-mineral}{高岭石}，
\qty{492.3}{kg}。钙\omterm{def:b1:rock-forming-minerals:groups}{长石}：\qty{3594.5}{mol}，\qty{928.0}{kg}。

\textbf{9.} 钠\omterm{def:b1:rock-forming-minerals:groups}{长石}：$2\times3813.6 = \qty{7627.2}{mol}$ 的 \ce{SiO2}，
\qty{458.2}{kg}，以及 \qty{3813.6}{mol} 的钠，\qty{87.7}{kg}。钙\omterm{def:b1:rock-forming-minerals:groups}{长石}：
无二氧化硅，以及 \qty{3594.5}{mol} 的钙，\qty{144.1}{kg}。

\textbf{10.} 钠\omterm{def:b1:rock-forming-minerals:groups}{长石}：$1/2.62 = \qty{0.382}{m^3}$；其\omterm{def:b1:weathering-clays:clay-mineral}{高岭石}，
$0.4923/2.6 = \qty{0.189}{m^3}$：固体损失其
体积的 \qty{50.4}{\percent}。钙\omterm{def:b1:rock-forming-minerals:groups}{长石}：$1/2.75 = \qty{0.364}{m^3}$；其\omterm{def:b1:weathering-clays:clay-mineral}{高岭石}，
$0.9280/2.6 = \qty{0.357}{m^3}$：仅损失 \qty{1.8}{\percent}。

\textbf{11.} 余下的空间变为孔隙，充以水和空气：富钠\omterm{def:b1:rock-forming-minerals:groups}{长石}花岗岩上的一种腐岩
很轻、多孔、疏松，其石英
颗粒散布于少量黏土中。钙\omterm{def:b1:rock-forming-minerals:groups}{长石}转化为体积几乎与其自身相当的
\omterm{def:b1:weathering-clays:clay-mineral}{高岭石}，因此风化后的辉长岩保留了大部分固体物质。

\textbf{12.} 钠（$z/r = 0.98$）和钙（2.0）为可溶阳离子，并
迁出；铝（5.6）为\omterm{def:b1:weathering-clays:hydrolysis}{水解}产物，全部留在\omterm{def:b1:weathering-clays:clay-mineral}{高岭石}中；
硅（10.0）处于边界，部分自钠\omterm{def:b1:rock-forming-minerals:groups}{长石}中迁出，完全不
自钙\omterm{def:b1:rock-forming-minerals:groups}{长石}中迁出，其硅可为\omterm{def:b1:weathering-clays:clay-mineral}{高岭石}所容纳。

\textbf{13.} 每个钠\omterm{def:b1:rock-forming-minerals:groups}{长石}一个 \ce{CO2}：\qty{3813.6}{mol}，$3813.6 \times 44.01 =
\qty{167.8}{kg}$.

\textbf{14.} 每个钙\omterm{def:b1:rock-forming-minerals:groups}{长石}两个：\qty{7189.0}{mol}，\qty{316.4}{kg}。

\textbf{15.} $7189.0 \times 61.02 = \qty{438.7}{kg}$ 的碳酸氢根。

\textbf{16.} 生物将其沉淀为碳酸钙，
\ce{Ca^2+ + 2 HCO3- -> CaCO3 + CO2 + H2O}：每个钙对应
两个 \ce{CO2} 中的一个返回空气，即 \qty{3594.5}{mol} 或 \qty{158.2}{kg}，而
另一个埋藏于碳酸盐中。

\textbf{17.} 钠在海底不形成碳酸盐：其碳酸氢根
不能以碳酸盐岩形式储存，长期储存量仅
依据钙和镁计算。

\textbf{18.} 一吨方解石溶解时消耗 \qty{9991}{mol} 的 \ce{CO2}，
沉淀时又全部释放：没有任何储存，而
只有钙镁硅酸盐的风化才会将二氧化碳
永久移出。

\textbf{19.} $\qty{200}{m^3.s^{-1}} \times \qty{0.3}{mol.m^{-3}} =
\qty{60}{mol.s^{-1}}$；一年（\qty{3.156e7}{s}）共 \qty{1.893e9}{mol}，或
$\num{1.893e9}/\num{50000} = \qty{37870}{mol}$ 每平方千米。

\textbf{20.} 每个钙\omterm{def:b1:rock-forming-minerals:groups}{长石}两个碳酸氢根：每年 \qty{9.467e8}{mol}，
$\num{9.467e8}\times\qty{278.20}{g} = \qty{263400}{t}$ 每年，或
每平方千米 \qty{5.27}{t}。

\textbf{21.} 所消耗的新鲜岩石为 $5.27/0.20 = \qty{26.3}{t}$，或
$\num{26340}/2700 = \qty{9.75}{m^3}$，即每平方千米每年：一层
每年厚 \qty{9.75}{\micro m}，每百万年 \qty{9.8}{m}。

\textbf{22.} 每个钙储存一个 \ce{CO2}：$\num{9.467e8} \times 44.01 =
\qty{4.17e10}{g}$，每年约 \qty{41700}{t}。

\textbf{23.} 这里只计入以溶液形式带走的物质。风化留下的\omterm{def:b1:weathering-clays:clay-mineral}{高岭石}
和石英还会被侵蚀作用带走，成为河流中的泥和沙，
它们均未计入；岩石中其他
\omterm{def:b1:rock-forming-minerals:mineral}{矿物}的风化亦未计入。若部分碳酸氢根来自钠或
钾\omterm{def:b1:rock-forming-minerals:groups}{长石}，它们每个化学式单位只给出一个碳酸氢根而非两个，
则须有更多岩石风化才能提供。

\textbf{24.} 由\cref{prop:b1:weathering-clays:co2}可知，每摩尔钙\omterm{def:b1:rock-forming-minerals:groups}{长石}
储存一摩尔 \ce{CO2}：$3594.5 \times \qty{44.01}{g} = \qty{158.19}{kg}$。
一吨钙\omterm{def:b1:rock-forming-minerals:groups}{长石}的风化可使
\textbf{158}~千克二氧化碳被永久移出。
\end{solution}
